\documentclass[a4paper,11pt]{article}

% 导入包
\usepackage[UTF8]{ctex}
\usepackage{geometry}
\usepackage{titlesec}
\usepackage{enumitem}
\usepackage[hidelinks]{hyperref}
\usepackage{fontawesome}
\usepackage{xcolor}
\usepackage{tabularx}

% 设置页面边距 - 严格控制
\geometry{left=0.7in,top=0.6in,right=0.7in,bottom=0.6in}

% 定义全局字体大小控制变量
\newcommand{\basefontsize}{\fontsize{12}{13}\selectfont}  % 基础正文大小，可直接修改磅数
\newcommand{\itemfontsize}{\fontsize{11}{12}\selectfont}   % 列表项字体大小
\newcommand{\descfontsize}{\fontsize{11}{12}\selectfont}   % 描述文字大小
\newcommand{\titlefontsize}{\fontsize{20}{19}\selectfont}  % 标题字体大小
\newcommand{\subsubtitlefontsize}{\fontsize{13}
{16}\selectfont} % 子标题字体大小
\newcommand{\subtitlefontsize}{\fontsize{15}
{12}\selectfont} % 子标题字体大小


% 定义苹果风格颜色
\definecolor{appleblue}{RGB}{0,0,30}      % 苹果蓝
\definecolor{applegray}{RGB}{99,99,102}      % 苹果灰
\definecolor{appledark}{RGB}{28,28,30}       % 苹果深色
\definecolor{textgray}{RGB}{85,85,85}        % 文本灰色

% 设置页面样式
\pagestyle{empty}

% URL样式设置
\urlstyle{same}

% 文档布局设置 - 严格间距控制
\raggedbottom
\raggedright
\setlength{\tabcolsep}{0in}
\setlength{\parindent}{0pt}
\setlength{\parskip}{0pt}

% 节标题格式设置 - 苹果风格简洁设计
\titleformat{\section}
    {\subtitlefontsize\bfseries\color{appleblue}}
    {}{0em}
    {}[\vspace{-10pt}{\color{applegray}\rule{\textwidth}{0.8pt}}]
\titlespacing{\section}{0pt}{-5pt}{-5pt}

% 自定义命令 - 严格间距控制
\newcommand{\resumeItem}[1]{
    \item\itemfontsize{#1}
}

\newcommand{\resumeSubheading}[4]{
    \vspace{3pt}
    \begin{tabular*}{\textwidth}{l@{\extracolsep{\fill}}r}
        \textbf{\subsubtitlefontsize\color{appledark}#1} \textit{\basefontsize\color{textgray} · #3} & {\basefontsize\color{applegray}#2} \\
    \end{tabular*}
    \vspace{2pt}
}

\newcommand{\resumeSubHeadingListStart}{\begin{itemize}[leftmargin=0pt,label={},itemsep=5pt,parsep=0pt,topsep=0pt]}
\newcommand{\resumeSubHeadingListEnd}{\end{itemize}\vspace{6pt}}
\newcommand{\resumeItemListStart}{\begin{itemize}[leftmargin=15pt,itemsep=2pt,parsep=0pt,topsep=3pt]}
\newcommand{\resumeItemListEnd}{\end{itemize}\vspace{4pt}}

\newcommand{\contactItem}[2]{
    \textcolor{textgray}{\basefontsize#1}~\textcolor{appleblue}{\basefontsize#2}\quad
}

\begin{document}

% 标题和联系信息 - 苹果风格简洁设计
\begin{center}
    {\titlefontsize\bfseries\color{appledark}  后端开发 - 刘龙昌\\     
    已转正成功，可长期实习六个月以上，周内到岗} \\
    \vspace{4pt}

    \vspace{12pt}
    \begin{tabular}{c}
        \contactItem{\faPhone}{18138256454}
        \contactItem{\faEnvelope}{liulongchang90@163.com}
        \contactItem{\faGithub}{\href{https://github.com/leochame}{github.com/leochame}}
        \contactItem{\faGlobe}{\href{https://leochame.github.io}{leochame.github.io}}\end{tabular}
\end{center}
\vspace{4pt}

% 教育经历
\section{教育经历}
\resumeSubHeadingListStart
    \item\resumeSubheading
    {杭州电子科技大学}{2022年08月 - 2026年06月}{计算机学与技术 · 本科}{}
\resumeSubHeadingListEnd
\vspace{-6pt}
% 实习经历
\section{实习经历}
\resumeSubHeadingListStart
    
    \item\resumeSubheading
    {\subtitlefontsize {腾讯}}{2025年04月 - 2025年09月}{广告营销服务线}
    {}{}
    \descfontsize{\color{textgray}{参与腾讯广告 \textbf{Agentic RAG} 核心服务 ChatBI 的系统设计与实现。 该服务旨在借助 LLM 完成用户自然语言到返回图表的实现，负责接收并解析用户的自然语言查询，通过调用 LLM 将其\textbf{工程化地转化为结构化的数据查询指令（Text2DSL）}，并调度下游数据系统的执行。}
\resumeItemListStart
    \resumeItem{为解决LLM知识库的源头数据质量瓶颈，设计并实现了一套\textbf{自动化需求文档评审系统。} 该系统利用LLM自动分析指标定义完整度，并将评审建议无缝回写至需求管理平台，构建了从“问题发现”到“协同修正”的工作流闭环，推动召回的正确率从50\% 到 80\%。}
    \resumeItem{为保障评审流程的系统健壮性，以\textbf{状态机模型}为核心设计了任务调度系统，并通过阻塞队列实现各环节异步解耦。在增强架构的可扩展性的同时，成功实现了复杂任务精确追踪与故障续跑。}
\resumeItemListEnd
    \item\resumeSubheading
    {iMile国际物流公司}{2024年09月 - 2025年02月}
    {}{}
    {\normalsize \textcolor{darkgray}{负责业务数据服务平台的后端开发，为前端可视化及运营分析提供API支持。在7个版本迭代中，主导开发了超20项数据处理与查询功能，支撑了20+国家、万级用户的数据需求。}\vspace{3pt}
    \resumeItemListStart
        \resumeItem{针对目标订单数据，基于“本地消息表”模式设计并实现了向下游3个系统的数据广播方案，确保了99.9\%的数据同步成功率与最终一致性，实现了跨系统的零失误传递。}
        \resumeItem{针对亿级数据查询接口10s+的响应瓶颈，主导全链路性能分析并引入MapReduce并行计算。最终将\textbf{接口响应时间从12s+降至3s+}，吞吐量提升6倍，稳定支撑50+业务方的高并发访问。}
        \resumeItem{为提升团队开发效率，\textbf{基于Java反射机制封装通用深拷贝工具}，解决了复杂对象深拷贝繁琐的痛点。不仅减少了50\%的业务样板代码，还落实技术文档，推广至多个业务线使用。}
    \resumeItemListEnd
\resumeSubHeadingListEnd
% 项目经验
\section{项目经验}
\resumeSubHeadingListStart
    \item\resumeSubheading
    {学院竞赛智慧社区平台}{2022年 - 2025年}
    {主要开发者 · 学院项目}{已上线运行3年}
    {\descfontsize\color{textgray}学院竞赛平台作为支撑学院学科竞赛的数字化基座，\textbf{基于冷热分离解决订量持续增长带来的在线存储性能瓶颈}，集成多级流水线传输技术，解决大文件传输抖动造成文件重传问题。服务学院90\%竞赛活动。}
    \resumeItemListStart
        \resumeItem{针对大文件传输因网络抖动等原因频繁触发重传，\textbf{采用分片+异步编程配合线程池}，将GB级别文件的上传时间压缩至3s内，重传率从25\%降至3\%}
        \resumeItem{设计实现\textbf{分布式定时任务系统}，通过Redis 的 ZSET+LIST双队列结构和数据库持久化方案，最高支撑日均1万+任务的精准定时发布}
    \resumeItemListEnd
\resumeSubHeadingListEnd
% 专业技能
\section{专业技能}
\vspace{0pt}
\begin{itemize}[leftmargin=15pt,itemsep=3pt,parsep=0pt,topsep=4pt]
\vspace{6pt}
    \resumeItem\resumeItem{
    \textbf{LLM与AI工程}：深入理解 Transformer 架构与注意力机制，熟悉 CoT、ReAct 等主流 Agent 设计思想，具备从0到1独立开发 多智能体协同的 LLM应用的工程实践能力。
}
    \resumeItem{\textbf{实践能力}：熟练SQL语句优化，MySQL性能调优，JAVA线上故障排查，以及线上问题定位}
    \resumeItem{\textbf{计算机基础}：熟悉页表、Trap、System Call、线程切换、TCP/UDP协议、拥塞控制等核心机制}
    \resumeItem{\textbf{Java基础}：熟悉GC（G1/ZGC）、类加载及内存模型，具备JVM调优及性能监控实战经验}
\end{itemize}
\end{document}